%! Properties of Real Numbers, Equality, and Inequality — Unit 1, Lesson 1.0 — Warm-Up (Answer Key)
\documentclass[10pt]{article}
\usepackage{algebra2-article}
\usepackage{algebra2-key}   % pulls in -boxes; do NOT also load -boxes
\newcommand{\TallMath}[1]{$\displaystyle #1\rule[-1.4em]{0pt}{3.2em}$}
\newcommand{\ck}{\ans{$\checkmark$}}

\begin{document}

\pageheader{Unit 1, Lesson 1.0}{Warm-Up --- Answer Key}

\vspace{0.08in}

\begin{notesbox}{Getting Ready: Skills We'll Use Today}
\small Welcome back. These three items are all Algebra 1 --- you can already do every one of them.
Today's lesson gives names to what you are about to do.

\vspace{4pt}
\textbf{1. Where does each number live?} Put a check in \emph{every} box that the number belongs to.
Several numbers belong to more than one.
\begin{center}
\renewcommand{\arraystretch}{1.45}
\begin{tabular}{|l|c|c|c|c|c|}
\hline
\textbf{Number} & \textbf{Natural} & \textbf{Whole} & \textbf{Integer} & \textbf{Rational} & \textbf{Irrational} \\
\hline
$-7$ & & & \ck & \ck & \\ \hline
$0$ & & \ck & \ck & \ck & \\ \hline
$\tfrac34$ & & & & \ck & \\ \hline
$\sqrt{9}$ & \ck & \ck & \ck & \ck & \\ \hline
$\sqrt{5}$ & & & & & \ck \\ \hline
$\pi$ & & & & & \ck \\ \hline
\end{tabular}
\end{center}
Two of these six numbers caused an argument at your table. Which two, and why?
\ansline{$\sqrt{9}$ and $0$. $\sqrt{9}=3$, so it is a \emph{natural number} --- the radical sign
does not make a number irrational; the \emph{value} decides. And $0$ is whole and an integer but
\emph{not} natural, since counting starts at $1$.}

\vspace{4pt}
\textbf{2. Fill in the blank --- these should take you no time at all.}
\begin{center}
\renewcommand{\arraystretch}{1.5}
\begin{tabular}{@{}l l l l@{}}
(a) \; $5+0=$ \ans{5} &
(b) \; $7\cdot 1=$ \ans{7} &
(c) \; $4+(-4)=$ \ans{0} &
(d) \; $2\cdot\tfrac12=$ \ans{1} \\
\end{tabular}
\end{center}
Every one of those four is an example of a rule that has a \emph{name}. Do you know any of the four
names? Write down any you know --- and it is completely fine to write ``none.''
\ansline{Accept anything, including ``none.'' The four names, handed out in Section 2 today:
(a) additive identity, (b) multiplicative identity, (c) additive inverse, (d) multiplicative
inverse (reciprocal).}

\vspace{4pt}
\textbf{3. Solve each one. Show every step.}
\begin{center}
\begin{minipage}[t]{0.45\linewidth}
\centering $2x-7=11$
\vspace{2pt}

\ans{$2x = 18$}

\ans{$x = 9$}
\vspace{0.55cm}
\end{minipage}
\begin{minipage}[t]{0.45\linewidth}
\centering $-3x>12$
\vspace{2pt}

\ans{$x < -4$}

\ans{(divided by $-3$, so the symbol reversed)}
\vspace{0.55cm}
\end{minipage}
\end{center}
One of those two problems made you do something \emph{extra} that the other one did not. Which one,
and what was the extra thing?
\ansline{The inequality $-3x>12$. Dividing both sides by $-3$ --- a \emph{negative} number ---
forced the $>$ to turn around into $<$. Nothing like that happened in the equation.}

\end{notesbox}

\end{document}
